
\documentclass{article}
\usepackage{colortbl}
\usepackage{makecell}
\usepackage{multirow}
\usepackage{supertabular}

\begin{document}

\newcounter{utterance}

\centering \large Interaction Transcript for game `cladder', experiment `full\_v1.5\_default', episode 4386 with Qwen3.5{-}4B.
\vspace{24pt}

{ \footnotesize  \setcounter{utterance}{1}
\setlength{\tabcolsep}{0pt}
\begin{supertabular}{c@{$\;$}|p{.15\linewidth}@{}p{.15\linewidth}p{.15\linewidth}p{.15\linewidth}p{.15\linewidth}p{.15\linewidth}}
    \# & \multicolumn{2}{c}{Player} && \multicolumn{2}{c}{Game Master} \\
    \hline

    \theutterance \stepcounter{utterance}  
    & & & \multicolumn{4}{p{0.6\linewidth}}{
        \cellcolor[rgb]{0.9,0.9,0.9}{
            \makecell[{{p{\linewidth}}}]{
                \texttt{\tiny{[P1$\langle$GM]}}
                \texttt{You are an expert in causal inference. The following question is not a typical commonsense query, but rather a meticulously designed question created by a professor specializing in causal inference, intended to assess the students' mastery of the course content.} \\
\texttt{Imagine a self{-}contained, hypothetical world with only the following conditions, and without any unmentioned factors or causal relationships: Gender has a direct effect on smoking and peanut allergy. Smoking has a direct effect on tar deposit. Tar deposit has a direct effect on peanut allergy. Gender is unobserved. For nonsmokers and with no tar deposit, the probability of being allergic to peanuts is 35\%. For nonsmokers and with high tar deposit, the probability of being allergic to peanuts is 63\%. For smokers and with no tar deposit, the probability of being allergic to peanuts is 25\%. For smokers and with high tar deposit, the probability of being allergic to peanuts is 49\%. For nonsmokers, the probability of high tar deposit is 22\%. For smokers, the probability of high tar deposit is 73\%. For smokers, would it be more likely to see being allergic to peanuts if the person had been a nonsmoker?} \\
\texttt{Start your answer with 'Yes' or 'No', followed by additional reasoning or evidence to support your explanation.} \\
            }
        }
    }
     \\ \\

    \theutterance \stepcounter{utterance}  
    & \multicolumn{4}{p{0.6\linewidth}}{
        \cellcolor[rgb]{0.561,0.737,0.561}{
            \makecell[{{p{\linewidth}}}]{
                \texttt{\tiny{[P1$\rangle$GM]}}
                \texttt{No} \\
\\ 
\texttt{**Reasoning:**} \\
\\ 
\texttt{To determine whether it would be more likely to see a person allergic to peanuts if they had been a nonsmoker rather than a smoker, we must compare the **counterfactual probability** of peanut allergy for a smoker under the condition that they were a nonsmoker (\$P(\textbackslash{}text\{Allergy\} | \textbackslash{}text\{Smoker\} \textbackslash{}to \textbackslash{}text\{Nonsmoker\})\$) against the **observed probability** of peanut allergy for a smoker (\$P(\textbackslash{}text\{Allergy\} | \textbackslash{}text\{Smoker\})\$).} \\
\\ 
\texttt{**Step 1: Calculate the observed probability of peanut allergy for smokers.**} \\
\texttt{We are given the conditional probabilities of high tar deposit for smokers:} \\
\texttt{*   \$P(\textbackslash{}text\{High Tar\} | \textbackslash{}text\{Smoker\}) = 0.73\$} \\
\texttt{*   \$P(\textbackslash{}text\{No Tar\} | \textbackslash{}text\{Smoker\}) = 1 {-} 0.73 = 0.27\$} \\
\\ 
\texttt{We are given the probabilities of peanut allergy based on smoking status and tar deposit:} \\
\texttt{*   \$P(\textbackslash{}text\{Allergy\} | \textbackslash{}text\{Smoker\}, \textbackslash{}text\{High Tar\}) = 0.49\$} \\
\texttt{*   \$P(\textbackslash{}text\{Allergy\} | \textbackslash{}text\{Smoker\}, \textbackslash{}text\{No Tar\}) = 0.25\$} \\
\\ 
\texttt{Using the law of total probability:} \\
\texttt{\$\$P(\textbackslash{}text\{Allergy\} | \textbackslash{}text\{Smoker\}) = P(\textbackslash{}text\{Allergy\} | \textbackslash{}text\{Smoker\}, \textbackslash{}text\{High Tar\}) \textbackslash{}times P(\textbackslash{}text\{High Tar\} | \textbackslash{}text\{Smoker\}) + P(\textbackslash{}text\{Allergy\} | \textbackslash{}text\{Smoker\}, \textbackslash{}text\{No Tar\}) \textbackslash{}times P(\textbackslash{}text\{No Tar\} | \textbackslash{}text\{Smoker\})\$\$} \\
\texttt{\$\$P(\textbackslash{}text\{Allergy\} | \textbackslash{}text\{Smoker\}) = (0.49 \textbackslash{}times 0.73) + (0.25 \textbackslash{}times 0.27)\$\$} \\
\texttt{\$\$P(\textbackslash{}text\{Allergy\} | \textbackslash{}text\{Smoker\}) = 0.3577 + 0.0675 = 0.4252\$\$} \\
\\ 
\texttt{**Step 2: Calculate the counterfactual probability of peanut allergy for smokers if they were nonsmokers.**} \\
\texttt{We need to evaluate \$P(\textbackslash{}text\{Allergy\} | \textbackslash{}text\{Smoker\} \textbackslash{}to \textbackslash{}text\{Nonsmoker\})\$. This requires the distribution of tar deposits for nonsmokers and the resulting allergy probabilities.} \\
\\ 
\texttt{We are given the probabilities of high tar deposit for nonsmokers:} \\
\texttt{*   \$P(\textbackslash{}text\{High Tar\} | \textbackslash{}text\{Nonsmoker\}) = 0.22\$} \\
\texttt{*   \$P(\textbackslash{}text\{No Tar\} | \textbackslash{}text\{} \\
            }
        }
    }
    & & \\ \\

    \theutterance \stepcounter{utterance}  
    & & & \multicolumn{2}{p{0.3\linewidth}}{
        \cellcolor[rgb]{0.9,0.9,0.9}{
            \makecell[{{p{\linewidth}}}]{
                \texttt{\tiny{[GM$|$GM]}}
                \texttt{No} \\
            }
        }
    }
    & & \\ \\

    \theutterance \stepcounter{utterance}  
    & & & \multicolumn{2}{p{0.3\linewidth}}{
        \cellcolor[rgb]{0.9,0.9,0.9}{
            \makecell[{{p{\linewidth}}}]{
                \texttt{\tiny{[GM$|$GM]}}
                \texttt{no} \\
            }
        }
    }
    & & \\ \\

    \theutterance \stepcounter{utterance}  
    & & & \multicolumn{2}{p{0.3\linewidth}}{
        \cellcolor[rgb]{0.9,0.9,0.9}{
            \makecell[{{p{\linewidth}}}]{
                \texttt{\tiny{[GM$|$GM]}}
                \texttt{game\_result = WIN} \\
            }
        }
    }
    & & \\ \\

\end{supertabular}
}

\end{document}
